\subsection{System Architecture}

Describe the three components (LLM, validator, reprompting pipeline) and how they interact. Include your flowchart (Figure A1). Explain design choices, e.g. why Docker-hosted W3C rather than the public API.


\subsection{Prompt Dataset}

Describe how you constructed your prompt set. How many prompts? How did you decide which ones to include? Did you deliberately design some to provoke semantic tag errors (e.g. tasks that structurally require \verb|<nav>|, \verb|<article>|, etc.)? This directly addresses the supervisor's feedback about the random prompt selection being unclear.

\subsection{Models \& Configurations}

List all models tested. Describe the parameter ranges (size, quantisation). Explain how you vary the maximum iteration count and how you determine when you've found the "optimal" cut-off — frame this as a quality-vs-cost tradeoff that you measure empirically.

\subsection{Validation Strategy}

This is the most critical methodological section. Address the supervisor's concern directly: can the W3C validator actually detect semantic misuse of <div>? If it cannot (likely the case — W3C validates syntax, not semantics), you need to be transparent about this limitation and introduce a complementary evaluation. Options to discuss:

\begin{itemize}
	\item A human evaluation rubric (e.g. annotators scoring semantic correctness)
	\item A secondary automated check (e.g. counting occurrences of semantic tags vs. <div> before and after)
	\item Possibly a separate LLM-as-judge evaluation
\end{itemize}

Clarify what "comparable output quality" means for RQ2 — define your quality metric explicitly.


\subsection{Cost Metric}

Define how you operationalise "cost" — API pricing per token, latency, or total tokens consumed per run. This addresses the supervisor's question about RQ2's cost dimension.